Lemma 1: For groups of order $p^2 q$, there must exist either a normal Sylow $p$-group or a normal Sylow $q$-group.

Lemma 2: If $K$ is a cyclic $p$-group, two homomorphisms $f, g : K → \text{Aut}(H)$ define
isomorphic semi-direct products $H \rtimes K$ if they have the same image.

Lemma 3: Semi-direct products $C_{q^n} ⋊ C_{p^m}$ ($q$ odd prime, $p ≠ q$) are classified up to
isomorphism by $r ∈ \{0, …, \min(m, d)\}$ where $d = v_p(q - 1)$, giving $\min(m, d) + 1$ classes. Every action $f$ belongs to exactly one class.

Lemma 4: If $|K| = p$, then two homomorphisms $f, g : K → \text{Aut}(H)$ define isomorphic semi-direct products $H \rtimes K$ if their images are conjugate subgroups of $\text{Aut}(H)$.

$|N| = p^2$ and by Schur-zassenhaus theorem there exists a complementary subgroup $U$ of $G$ with $|U| = q$ such that $G \cong N \rtimes U$. We have $U \cong C_q$.

## Case 1.1: $N \cong C_{p^2}$
$\text{Aut}(N) = C_{p(p-1)}$
Let $\phi : C_q \rightarrow C_{p(p-1)}$. Done by Lemma 3 when $p$ is odd. When $p = 2$, we get $\phi : C_q \rightarrow C_2$ so only trivial homomorphism is possible as $q \nmid 2$ for odd $q$.

(For $p = 2$, only trivial homomorphism)
(For $q = 2$, then we get 2 homomorphisms)

## Case 1.2: $N \cong C_p \times C_p$
$\text{Aut}(N) \cong \text{GL}_2(p)$
Let $\phi : C_q \rightarrow \text{GL}_2(p)$
### Special case $p = 2$
$GL_2(2) \cong S_3$
Either $\phi$ is a trivial action, or its image is a subgroup of order $q$ in $S_3$, hence $q = 3$. Only one subgroup of order $3$ so gives one group.

### Special case $q = 2$
$\text{Im}(\phi)$ is trivial or $\langle A \rangle$ where $A \in GL_2(p)$ is of order $2$. By lemma 4 we only need to look at similarity classes and pick a representative from each class. Then $A^2 = 1$ implies minimal polynomial of $A$ divides $x^2 - 1$ and since $p$ is odd, it has distinct factors hence diagonalisable. Hence each such matrix similar to a diagonal matrix $\text{diag}(\alpha, \beta)$ such that $\alpha^2 = 1, \beta^2 = 1$ hence $\alpha = \pm 1, \beta = \pm 1$ ($t^2 - 1$ has at most $2$ roots because we are over a field). So $\alpha = \beta = 1$ is order $1$ so we don't consider it. We can check there are two similarity classes $\alpha = 1, \beta = -1$ (or switch them) and one for $\alpha = \beta = -1$.

Side note: Lemma 4 not needed because the general condition $\phi_2(u) = \alpha \circ \phi_1(\beta(u)) \circ \alpha^{-1}$ becomes equivalent to similarity because $\text{Aut}(C_q) = C_1$ as $q = 2$ so $\beta = \text{id}$ and that $\phi_1, \phi_2$ are determined by where they send the generator as $U$ is cyclic.

# Case 2: There exists a normal Sylow $q$-group

$|N| = q$ and by Schur-zassenhaus theorem there exists a complementary subgroup $U$ of $G$ with $|U| = p^2$ such that $G \cong N \rtimes U$. We have $N \cong C_q$. $\text{Aut}(N) \cong C_{q-1}$.

## Case 2.1 $U \cong C_{p^2}$
Let $\phi : C_{p^2} \rightarrow C_{q-1}$. Done by Lemma 3 when $q$ is odd. Only trivial homomorphism possible when $q = 2$.

(When $p = 2$, have to see if $q$ is $1$ mod $4$)
(When $q = 2$, we get only trivial homomorphism)

## Case 2.2 $U \cong C_p \times C_p$
Let $\phi : C_p \times C_p \rightarrow C_{q-1}$

$\text{Im}(\phi) \cong (C_p \times C_p) / \ker(\phi)$ so order of image is either $1, p$ or $p^2$. If it is $1$, then we have the trivial homomorphism giving $G \cong C_p \times C_p \times C_{q}$. If it is $p^2$, then we must have $\ker(\phi) = 1$ but $\text{Im}(\phi)$ would then be a non-cyclic subgroup of the cyclic group $C_{q-1}$, impossible. If it is $p$, then ?????

(When $p = 2$, one trivial and one non-trivial?)
(When $q = 2$, we get only trivial homomorphism)
# p = 2, q > 3, q is 1 mod 4
1 (1.1) + 1 (1.2) + 3 (2.1) + 2 (2.2) = 6

According to formula: 5, we get $C_{4q}$ again from case 2.1 (trivial homomorphism)
$C_{4q}$ (1.1)
$C_2 \times C_2 \times C_q$ (1.2)
$C_q \rtimes_1 C_4$ (2.1)
$C_q \rtimes_2 C_4$ (2.1)
$C_q \rtimes (C_2 \times C_2)$ (2.2)
# p = 2, q > 3, q is 3 mod 4
1 (1.1) + 1 (1.2) + 2 (2.1) + 2 (2.2) = 6

$C_{4q}$ (1.1)
$C_2 \times C_2 \times C_q$ (1.2)
$C_q \rtimes C_4$ (2.1)
$C_q \rtimes (C_2 \times C_2)$ (2.2)

According to formula: 4, we get $C_{p^2 q}$ again from case 2.1 (trivial homomorphism)
# p = 2, q = 3
1 (1.1) + 2 (1.2) + 2 (2.1) + 2 (2.2) = 7
There are only 5 so some 2 are isomorphic here

$C_{12}$ (1.1)
$C_2 \times C_2 \times C_3$ (1.2)
$(C_2 \times C_2) \rtimes C_3$ (1.2)
$C_3 \rtimes C_4$ (2.1)
$C_3 \rtimes (C_2 \times C_2)$ (2.2)
# q = 2
2 (1.1), 2 (1.2), 1 (2.1, $C_{2p^2}$ repeated), 1 (2.2) total = 5